\begin{table}[h]
\centering
\caption{Falsification Test: Effect of Exposure to Homeowners on Highway Placement}
\label{tab:robustness/falsification_homeowners}
\begin{threeparttable}
\begin{tabular*}{0.6\textwidth}{@{\extracolsep{\fill}}l*{1}{r}}
\toprule
 & Coefficient \\
\midrule
Residential & \makecell[tr]{-0.875{**} \\ (0.338)} \\
Exposure to Homeowners & \makecell[tr]{-0.027 \\ (0.262)} \\
Residential x Exposure to Homeowners & \makecell[tr]{0.033 \\ (0.267)} \\
Est. Highway Suitability & \makecell[tr]{1.339{***} \\ (0.145)} \\
R-squared & 0.039 \\
Observations & 5236 \\
\bottomrule
\end{tabular*}
\begin{tablenotes}[flushleft]
\footnotesize
\item Estimates from a Poisson pseudo-maximum likelihood model (PPML) regression of an indicator for highway construction from 1940-1959 on the covariates listed, as well as controls for housing, geographic, and demographic characteristics. The sample is restricted to squares that are outside the highway corridor. Residential is a binary variable indicating whether the grid square is zoned for residential use. Exposure to Homeowners is a location-specific weighted average of the share of homeowners in each square. Weights are exponential decay function of the straight-line distance between squares and set to zero beyond a distance of 3,000 meters. This measure is normalized within each city to account for city-level differences. Standard errors, reported in parenthesis, are \textcite{conley_gmm_1999} spatial HAC standard errors with a 1,000-meter distance cutoff. *** p < 0.01, ** p < 0.05, * p < 0.1
\end{tablenotes}
\end{threeparttable}
\end{table}
